\section{Berücksichtigung der Normalkraft und der Bohrung}
\subsection{Fläche und FTM des Profils}
Die Bohrung wird Material in dem gefährlichem Querschnitt abtrennen. Die Bohrung ist symmetrisch um die y-Achse, und geht durch eine Trapezförmige Querschnitt des Profils.

\begin{figure}[H]
  \centering
  \begin{tikzpicture}
    \draw[blue] (2.5, 2) -- (2.5, 4);
    \fill[pattern=horizontal lines, pattern color=blue] (2.5, 2) -- (3, 4) -- (2.5, 4);

    \draw[blue] (7.5, 2) -- (7, 0);
    \fill[pattern=horizontal lines, pattern color=blue] (7.5, 2) -- (7, 0) -- (7.5, 0);

    \draw[ultra thick] (0,0) -- (2,0) -- (3, 4) -- (0, 4) -- (0,0);
    \draw[ultra thick] (5,0) -- (7.5,0) -- (7.5, 4) -- (5, 4) -- (5,0);
  \end{tikzpicture}
  \caption{Vereinfachung des Profils}
  \label{fig:trapezinrechteck}
\end{figure}
Wenn die Biegung um die z-Achse passiert, kann das Querschnitt als im \figref{fig:trapezinrechteck} zu sehen ist vereinfacht werden. Dieser Vereinfachung ändert $I_z$ nicht.

Die Berechung des Flächenträgheitmoments den gebohrten Querschnitte ist so einfach:
\begin{align*}
  &I_z = 2 \cdot \frac{d^3 \cdot t}{12} = 138666,663 \, mm^4
\end{align*}

Die neue Fläche und Schwerpunkt muss auch berechnet werden:
\begin{align*}
  &A = A_{U240} - A_{bohr} = 4230 - 2 \cdot 40 \cdot 13 = 3190 \, mm^2 \\
  &s
\end{align*}
